(1) 在阱内设
\[
k=\sqrt{2m(E+V_0)}/\hbar>0,\qquad
\kappa=\sqrt{-2mE}/\hbar>0,
\]
满足 \(k^2+\kappa^2=2mV_0/\hbar^2\)。偶宇称解为阱内 \(A\cos kx\)，阱外 \(A\cos(ka)e^{-\kappa(|x|-a)}\)；在 \(x=a\) 处导数连续给
\[
k\tan(ka)=\kappa\qquad\text{（偶态）}.
\]
奇宇称解为阱内 \(A\sin kx\)，阱外 \(A\,\operatorname{sgn}(x)\sin(ka)e^{-\kappa(|x|-a)}\)；导数连续给
\[
-k\cot(ka)=\kappa\qquad\text{（奇态）}.
\]

(2) 令 \(z=ka\in(0,\lambda)\)，则 \(\kappa a=\sqrt{\lambda^2-z^2}\)。偶态分支依次位于 \(z\in(0,\pi/2),(\pi,3\pi/2),\ldots\)；奇态分支依次位于 \((\pi/2,\pi),(3\pi/2,2\pi),\ldots\)。每跨过 \(\lambda=j\pi/2\)（\(j=1,2,\ldots\)）就允许出现一个新束缚分支，严格在阈值的态 \(E=0\) 不平方可积。因此总束缚态数为
\[
N_{\rm bound}=\left\lceil\frac{2\lambda}{\pi}\right\rceil
\]
（\(\lambda>0\)），包括始终存在的偶宇称基态。

(3) 对偶态，在 \(-\infty\) 到 \(\infty\) 积分 \(|\psi|^2\)，得到
\[
A_{\rm even}^{-2}=a+\frac{\sin(2ka)}{2k}+\frac{\cos^2(ka)}{\kappa}.
\]
对奇态则有
\[
A_{\rm odd}^{-2}=a-\frac{\sin(2ka)}{2k}+\frac{\sin^2(ka)}{\kappa}.
\]
两式皆严格正。对于一维连续有限势的平方可积本征态，同一能量的两个独立解的 Wronskian 必为常数，且在无穷远衰减使其为零，因此两解线性相关，说明束缚能级不简并。

(4) 在 \(V_0\to\infty\) 且固定级次 \(n=1,2,\ldots\) 时 \(\kappa a\gg1\)，边界趋近 Dirichlet 条件 \(x=\pm a\)，故
\[
E_n=-V_0+\frac{\hbar^2\pi^2n^2}{8ma^2}+o(1)
\]
（对应一条宽 \(2a\) 的无限深方阱的低级能量）。在浅阱 \(V_0\to0^+\) 时仅一个偶态；由 \(ka\ll1\) 与 \(k\tan(ka)\simeq k^2a\)，以及 \(k^2\simeq2mV_0/\hbar^2\)，得到 \(\kappa\simeq2mV_0a/\hbar^2\)，因此
\[
E_{\rm ground}\simeq-\frac{2mV_0^2a^2}{\hbar^2}.
\]
这一束缚态能量在浅阱极限接近连续谱边缘，但只要 \(V_0>0\) 就仍为负。
